% =============================================================================
% Section 9: Observed Results
% =============================================================================
% AGENT: This section is ONLY filled after unblinding (Wave 6).
% Do NOT populate during Wave 5. All placeholders below are for post-unblinding use.
% =============================================================================

\section{Observed Results}
\label{sec:observed}

% AGENT: This entire section is completed AFTER the unblinding procedure
% described in Appendix~\ref{sec:blinding}. Do NOT fill any subsection
% until blinding_status in STATE.md is set to UNBLINDED.

\subsection{Observed Upper Limit}
\label{sec:observed:limit}

% AGENT: Fill observed limit from Fitter.observed_limit() result
% Include observed value and comparison with expected bands.

The observed 95\% CL$_\text{s}$ upper limit on the signal strength
parameter $\mu$ is:
\begin{equation}
  \mu_\text{obs}^{95} = [\text{value}]
\end{equation}

% AGENT: Fill observed Brazil band plot from Fitter output with observed data.

\begin{figure}[htbp]
  \centering
  % AGENT: Replace with actual observed Brazil band plot path
  % \includegraphics[width=0.7\textwidth]{figures/brazil_band_observed.pdf}
  \caption{Observed and expected 95\% CL$_\text{s}$ upper limits on the
           signal strength $\mu$. The observed limit is shown as a solid
           line, with the expected median and $\pm 1\sigma$, $\pm 2\sigma$
           bands.}
  \label{fig:brazil_observed}
\end{figure}

\begin{table}[htbp]
  \centering
  \caption{Observed and expected 95\% CL$_\text{s}$ upper limits on $\mu$.}
  \label{tab:observed_limit}
  \begin{tabular}{lr}
    \toprule
    Quantity & Value \\
    \midrule
    % AGENT: Fill from Fitter output after unblinding
    Observed            & [value] \\
    $-2\sigma$ expected & [value] \\
    $-1\sigma$ expected & [value] \\
    Median expected     & [value] \\
    $+1\sigma$ expected & [value] \\
    $+2\sigma$ expected & [value] \\
    \bottomrule
  \end{tabular}
\end{table}

\subsection{Post-fit Yields}
\label{sec:observed:yields}

% AGENT: Fill post-fit yields from Fitter.fit() result with observed data.
% Show comparison of pre-fit and post-fit yields in all regions.

\begin{table}[htbp]
  \centering
  \caption{Post-fit yields in all analysis regions after the fit to observed data.}
  \label{tab:postfit_yields}
  \begin{tabular}{lrrr}
    \toprule
    Process & CR (post-fit) & VR (post-fit) & SR (post-fit) \\
    \midrule
    % AGENT: Fill from Fitter post-fit results
    Background 1  & -- & -- & -- \\
    Background 2  & -- & -- & -- \\
    Background 3  & -- & -- & -- \\
    \midrule
    Total background & -- & -- & -- \\
    Data             & -- & -- & -- \\
    \bottomrule
  \end{tabular}
\end{table}

\subsection{Post-fit Diagnostics}
\label{sec:observed:diagnostics}

% AGENT: Fill post-fit diagnostics from Diagnostics.run_all_diagnostics()
% with observed data. Compare with Asimov diagnostics from Section 8.

\subsubsection{Post-fit Nuisance Parameter Pulls}

% AGENT: Include post-fit NP pull plot with observed data.
% Flag any pulls exceeding +/-2 sigma for investigation.

\begin{figure}[htbp]
  \centering
  % AGENT: Replace with actual post-fit pull plot from observed data fit
  % \includegraphics[width=0.8\textwidth]{figures/np_pulls_observed.pdf}
  \caption{Post-fit nuisance parameter pulls from the fit to observed data.}
  \label{fig:np_pulls_observed}
\end{figure}

\subsubsection{Post-fit Correlation Matrix}

\begin{figure}[htbp]
  \centering
  % AGENT: Replace with actual post-fit correlation matrix from observed data
  % \includegraphics[width=0.7\textwidth]{figures/correlation_matrix_observed.pdf}
  \caption{Correlation matrix from the fit to observed data.}
  \label{fig:correlation_observed}
\end{figure}

\subsection{Comparison with Expected Results}
\label{sec:observed:comparison}

% AGENT: Compare observed results with pre-unblinding expectations.
% Document any significant deviations and their interpretation.

The observed limit of $\mu < $ [value] is [consistent with / differs from]
the expected limit of $\mu < $ [expected value]. The ratio of observed to
expected limits is [ratio], corresponding to a deviation of approximately
[N]$\sigma$ from the background-only expectation.

\subsection{Post-Unblinding Problems}
\label{sec:observed:problems}

% AGENT: If any Category A/B/C problems were classified after unblinding,
% document them here. Read from STATE.md post_unblinding_problems.
% Category A: cosmetic (proceed), Category B: requires re-blinding,
% Category C: requires external review.
% If no problems: state "No post-unblinding problems were identified."

No post-unblinding problems were identified.
% AGENT: Update the above statement based on STATE.md post_unblinding_problems
% field. If problems exist, describe each with its category, description,
% resolution, and impact on the result.
